% =====================================================================
%  PREÂMBULO — formatação. Em geral não é preciso mexer aqui.
% =====================================================================

% ---------- Idioma e codificação ----------
\usepackage[T1]{fontenc}
\usepackage[utf8]{inputenc}
\IfFileExists{brazilian.ldf}{\usepackage[english,brazilian]{babel}}{\usepackage[english]{babel}}

% ---------- Fonte Arial (equivalente livre: Helvetica) ----------
\usepackage[scaled=0.92]{helvet}
\renewcommand{\familydefault}{\sfdefault}
\usepackage{microtype}

% ---------- Página: A4, margens 3 cm (sup./esq.) e 2 cm (inf./dir.) ----------
\usepackage[a4paper,top=3cm,left=3cm,bottom=2cm,right=2cm]{geometry}

% ---------- Espaçamento 1,5 e recuo de 1,25 cm ----------
\usepackage{setspace}
\onehalfspacing
\setlength{\parindent}{1.25cm}
\setlength{\parskip}{0pt}
\usepackage{indentfirst}

% ---------- Pacotes gerais ----------
\usepackage{amsmath,amssymb}
\usepackage{siunitx}
\sisetup{output-decimal-marker={,}, group-separator={.}, per-mode=symbol}
\usepackage{graphicx}
\graphicspath{{figuras/}}
\usepackage{booktabs,tabularx,array,multirow}
\usepackage{float}
\usepackage[dvipsnames]{xcolor}
\usepackage{tikz}
\usetikzlibrary{positioning,arrows.meta,shapes.geometric,fit}
\usepackage{enumitem}
\setlist{nosep,leftmargin=1.25cm}
\usepackage{ifthen}
\newboolean{rascunho}

% ---------- Títulos: "1 INTRODUÇÃO" (12 pt, negrito, maiúsculas) ----------
\usepackage{titlesec}
\titleformat{\section}{\normalsize\bfseries}{\thesection}{0.5em}{\MakeUppercase}
\titleformat{\subsection}{\normalsize\bfseries}{\thesubsection}{0.5em}{}
\titleformat{\subsubsection}{\normalsize\bfseries\itshape}{\thesubsubsection}{0.5em}{}
\titlespacing*{\section}{0pt}{12pt}{6pt}
\titlespacing*{\subsection}{0pt}{12pt}{6pt}
\titlespacing*{\subsubsection}{0pt}{6pt}{6pt}
% Seções sem número (RESUMO, ABSTRACT, REFERÊNCIAS...)
\newcommand{\secaosemnumero}[1]{%
  \section*{#1}\addcontentsline{toc}{section}{#1}}

% ---------- Legendas: "Figura 1 — Título" (10 pt, centralizada) ----------
\usepackage{caption}
\captionsetup{font=small,labelsep=endash,justification=centering,singlelinecheck=false,skip=6pt}
\captionsetup[table]{position=top}
\DeclareCaptionLabelSeparator{travessao}{\ --\ }
\captionsetup{labelsep=travessao}
% Fonte abaixo da figura/tabela: \fonte{Autor (2026).}
\newcommand{\fonte}[1]{\par\vspace{4pt}{\centering\small\linespread{1}\selectfont Fonte: #1\par}}

% Novos ambientes flutuantes usados nos modelos: Quadro e Gráfico
\usepackage{newfloat}
\DeclareFloatingEnvironment[fileext=loq,listname={Lista de Quadros},name=Quadro,placement=htbp]{quadro}
\DeclareFloatingEnvironment[fileext=log2,listname={Lista de Gráficos},name=Gráfico,placement=htbp]{grafico}
\captionsetup[quadro]{position=top}

% Nomes em português independentemente do idioma carregado
\AtBeginDocument{%
  \renewcommand{\figurename}{Figura}%
  \renewcommand{\tablename}{Tabela}%
  \renewcommand{\abstractname}{Resumo}%
}

% ---------- Siglas (Dica 4: por extenso na 1.ª vez, depois só a sigla) ----------
\usepackage[printonlyused,nohyperlinks]{acronym}

% ---------- Links (precisa vir ANTES do abntex2cite) ----------
\usepackage[hidelinks]{hyperref}

% ---------- Citações ABNT: \citeonline{chave} -> Autor (2025); \cite{chave} -> (AUTOR, 2025) ----------
\IfFileExists{abntex2cite.sty}{%
  \usepackage[alf,abnt-emphasize=bf,abnt-etal-cite=2,abnt-etal-list=0,abnt-and-type=e]{abntex2cite}%
  \citebrackets()%
}{%
  \usepackage[round,authoryear]{natbib}
  \newcommand{\citeonline}{\citet}
  \bibliographystyle{plainnat}
}

% ---------- Orientações de escrita (somem quando rascunho = false) ----------
\newcommand{\escrever}[1]{%
  \ifthenelse{\boolean{rascunho}}{{\color{BrickRed}\itshape [Escrever: #1]}}{}}
\newcommand{\dica}[1]{%
  \ifthenelse{\boolean{rascunho}}{%
    \par\noindent\fcolorbox{BrickRed}{red!4}{\parbox{\dimexpr\linewidth-2\fboxsep-2\fboxrule}{\small\color{BrickRed}\textbf{Orientação:} #1}}\par}{}}
